\textbf{COLUNI 2022 -- Questão 26}\par\medskip Leia o texto a seguir.

\par

\textbf{Ostarina: a substância que causou afastamento de atletas dos Jogos Olímpicos de Tóquio}

\par

A ostarina ($C_{19}H_{14}F_3N_3O_3$) é um medicamento experimental pertencente à classe dos moduladores seletivos de receptores androgênicos (SARMs, sigla em inglês). Isso significa que a substância é capaz de se ligar às proteínas receptoras de androgênio (um hormônio masculino), causando o crescimento dos músculos. A substância e os produtos que a contenham são considerados ilegais pela agência regulatória \textit{Food and Drug Administration} (FDA) dos Estados Unidos, órgão regulador. No Brasil, a Agência Nacional de Vigilância Sanitária (Anvisa) publicou, em abril deste ano, a Resolução (RE) 791/2021, que proibiu a comercialização, a distribuição, a fabricação, a importação, a manipulação, a propaganda e o uso de produtos que contenham SARM, em função dos efeitos desconhecidos no corpo humano a longo prazo. Disponível em: https://g1.globo.com/ciencia-e-saude/noticia/2021/08/06/ostarina-entenda-o-que-e-a-substancia-que-causou-o-afastamento-de-tandara-dos-jogos-olimpicos.ghtml. Acesso em: 03 set. 2021 (Adaptado)

\par

Considere a equação química desta reação em que a ostarina está envolvida:

\par

$C_{19}H_{14}F_3N_3O_3 + H_2O \rightarrow C_8H_5F_3N_2 + C_{11}H_{11}NO_4$

\par

Analise as afirmativas abaixo, atribuindo V para a(s) verdadeira(s) e F para a(s) falsa(s):

\par

(   ) A equação química está corretamente balanceada.\par
(   ) A ostarina apresenta átomos de metais e ametais em sua composição.\par
(   ) Uma molécula de água apresenta 18 prótons.\par
(   ) A ostarina é um dos reagentes da reação.

\par

Assinale a sequência CORRETA:\begin{enumerate}
\item V, F, F, V.
\item F, F, V, V.
\item V, F, V, V.
\item F, V, F, F.\end{enumerate}
